\chapter{Introdução}
\label{cap:introducao}

As arboviroses transmitidas por \textit{Aedes aegypti} e \textit{Aedes albopictus}, em especial dengue (DENV), chikungunya (CHIKV) e Zika (ZIKV), estão entre as principais ameaças à saúde pública global \cite{gould2017}. Estima-se que 5,66 bilhões de pessoas vivam em áreas ambientalmente adequadas à transmissão de dengue, chikungunya e Zika, com subnotificação acentuada em regiões de baixa vigilância \cite{lim2025a}. Entre elas, a dengue é a de maior magnitude epidemiológica nas Américas, onde o número de casos cresceu de forma acentuada na última década, com epidemias recorrentes e expansão territorial da doença \cite{lim2025a,parums2025}. O impacto sanitário é elevado, com estimativas globais da ordem de 390 milhões de infecções anuais, além de morbidades crônicas e sobrecarga frequente dos sistemas de saúde \cite{parums2025}.

Fatores ecológicos e sociais se reforçam mutuamente. Mudanças climáticas alteram a distribuição dos vetores, favorecendo a expansão de áreas adequadas à transmissão. Sob cenários de aquecimento mais intenso, projeções indicam mais população exposta em diferentes regiões do mundo \cite{ryan2019}. Revisões sistemáticas mostram co-distribuição entre baixa escolaridade, precariedade habitacional, interrupções no abastecimento de água e maior vulnerabilidade à infecção em áreas urbanas densas \cite{power2022}. A dengue varia no espaço conforme o contexto ambiental e socioeconômico.

No Brasil, a circulação simultânea das três arboviroses caracteriza o período recente. Em 11 de novembro de 2015, o Ministério da Saúde declarou Emergência em Saúde Pública de Importância Nacional diante da alteração do padrão de ocorrência de microcefalias; a associação epidemiológica com a infecção congênita pelo vírus Zika foi descrita posteriormente. A série histórica do Ministério da Saúde, atualizada até a Semana Epidemiológica 47 de 2023, registra 213.350 casos prováveis de Zika em 2016 \cite{portaria1813,ms_zika}. Apenas em 2016, dengue, chikungunya e Zika geraram gastos da ordem de R\$ 2,3 bilhões \cite{2018-teich-arinelli-fahham}. A dengue concentrou o maior volume de casos e a dispersão territorial mais ampla, presente em quase todos os municípios do país \cite{pescarini2022}.

O Informe Semanal nº 18 do Ministério da Saúde, atualizado em 11 de junho de 2024, registrou 5.831.871 casos prováveis até a Semana Epidemiológica 23 e incidência de 2.872,0 por 100 mil habitantes; no mesmo recorte de 2023, a incidência era de 627,9 por 100 mil \cite{ms_informe2024}. Em 4 de dezembro de 2024, dados do painel ministerial reportados pela Agência Brasil indicavam 6.590.575 casos prováveis e incidência de 3.245 por 100 mil habitantes; esse valor é uma fotografia publicada naquela data, não o fechamento final do ano \cite{agenciabrasil_dengue2024}.

Muitos estudos avaliam os padrões de dengue e das demais arboviroses em escalas espaciais menores, e a cidade e o estado do Rio de Janeiro estão entre os territórios mais documentados no país. No Rio de Janeiro, estudos identificaram soroprevalência elevada para as três arboviroses e concentração de casos em áreas socialmente mais vulneráveis \cite{2023-souza-santos-sobral-perisse}. No município do Rio de Janeiro, a incidência das arboviroses concentra-se em espaços com piores indicadores socioeconômicos \cite{2021-queiroz-medronho}; na região Norte, a incidência de dengue acompanha vulnerabilidade social e urbanização \cite{2023-gomes-jesus-quaresma}. Essa heterogeneidade no espaço fundamenta a identificação de áreas de maior risco.

A heterogeneidade espacial e temporal orienta a estratificação de risco para a vigilância e o controle das arboviroses. Mapas de risco e detecção de agrupamentos espaciais fornecem informação territorial para o planejamento da vigilância e das ações de controle \cite{parums2025,liu2025}. Revisões da literatura apontam que a evidência disponível sobre a efetividade de diferentes intervenções de controle vetorial ainda é limitada \cite{bowman2016}.

Na literatura, a análise espacial da dengue e de outras arboviroses passou de abordagens descritivas para modelagens mais complexas. Estudos iniciais usaram índices de autocorrelação espacial, como o Índice de Moran \cite{moran1950,anselin1995}, para medir a dependência geográfica da doença em nível municipal e intramunicipal. \citeonline{almeida2009}, por exemplo, descreveram a concentração de casos na Zona Oeste do Rio de Janeiro, associada à falta de infraestrutura sanitária geral.

Em escalas distintas, \citeonline{rodrigues2018} e \citeonline{vernal2021} mapearam a variação espacial da incidência em comunidades vulneráveis e em municípios paulistas. Os resultados indicaram que indicadores estruturais básicos não bastavam para descrever a heterogeneidade da transmissão, que variava também com características sociodemográficas e geoclimáticas.

\citeonline{lowe2011} usaram modelos estatísticos baseados em anomalias climáticas sazonais para construir sistemas de alerta precoce. O efeito das temperaturas extremas e das chuvas na incidência e na dinâmica epidêmica foi analisado por meio de séries temporais e modelos espaço-temporais em diferentes regiões endêmicas \cite{xavier2021, ribeiro2025, alves2025}.

Em escala nacional, análises espaço-temporais associaram o risco de dengue a extremos hidrometeorológicos e à urbanização, incluindo mudanças na distribuição da doença em regiões do Sul e da Amazônia \cite{lee2021, lowe2021}. Modelos de difusão indicam maior adequação ambiental para mosquitos vetores em áreas temperadas da Europa \cite{barman2025}, enquanto a mortalidade segue concentrada em populações socialmente vulneráveis \cite{cardim2025}.

Para delimitar \textit{clusters} de alto risco, o método de Kulldorff tem sido usado para detectar \textit{clusters} espaço-temporais e rotas de disseminação de dengue, Zika e chikungunya, em contextos internacionais \cite{freitas2022} e no território nacional. No Brasil, \citeonline{palasio2023} investigaram \textit{clusters} espaço-temporais de Zika e chikungunya; \citeonline{almeida2023} identificaram \textit{clusters} simultâneos para as três arboviroses no estado do Rio de Janeiro, com vulnerabilidade persistente de regiões metropolitanas centrais; e revisões do cenário nacional tratam a dengue como desafio endêmico contínuo \cite{2020-magalhaes-et-al}.

Comparações entre as janelas circular e flexível indicam desempenho similar, com tendência da circular a delimitar áreas maiores e da flexível a delimitar áreas menores \cite{melo2022}. Ajustes no teto populacional refinam a sensibilidade da vigilância local \cite{canal2017}.

Esses métodos foram aplicados sobretudo em recortes territoriais específicos ou em períodos delimitados. Estudos sobre a expansão da dengue no Brasil descrevem ondas de difusão de regiões metropolitanas para cidades médias e pequenas, condicionadas por clima, população e conectividade \cite{barcellos2014}. Persiste, assim, uma lacuna de trabalhos que apliquem a varredura espacial de forma contínua a todo o território nacional por um período longo, com dados recentes, e que comparem o efeito de diferentes estratificações sobre a geografia dos excessos. O presente trabalho busca preencher essa lacuna ao cobrir os 5.564 municípios da malha de 2010 em 16 anos, de 2010 a 2025, incluindo o ano epidêmico de 2024, e ao contrapor oito formulações da mesma varredura, com e sem padronização indireta. O acompanhamento contínuo de toda a extensão do país viabiliza captar as dinâmicas de interiorização e a evolução temporal da doença, oferecendo subsídios para uma compreensão de conjunto do agravo.

Nesse sentido, o objetivo geral deste trabalho é descrever a distribuição espacial do risco de dengue no Brasil entre 2010 e 2025. Os objetivos específicos são descrever a distribuição da doença em conjunto com variáveis ambientais, climáticas e sociais; identificar e delimitar zonas de alto risco em todo o território nacional; e avaliar como a geografia dessas zonas se altera quando a varredura controla diferentes conjuntos de covariáveis e quando o ano epidêmico de 2024 é retirado da série. As comparações com as covariáveis são descritivas e não são interpretadas como efeitos causais. Para delimitar as zonas, emprega-se a janela circular clássica, cujas extensões devem ser interpretadas com cautela.

Além desta Introdução, o texto está organizado em três capítulos. O Capítulo~\ref{cap:matmet} apresenta a área de estudo, as bases de dados, o tratamento das informações, as medidas de autocorrelação espacial, a estatística de varredura espacial e a comparação entre as formulações, além dos recursos computacionais empregados. O Capítulo~\ref{cap:resultados} traz os resultados descritivos, a caracterização socioespacial e climática, a autocorrelação espacial e a análise dos \textit{clusters}, com a comparação entre as formulações e a sensibilidade à exclusão de 2024. O Capítulo~\ref{cap:discussao} discute os achados e reúne as conclusões, as limitações e as perspectivas. Os apêndices complementam o texto com as tabelas completas de \textit{clusters} e as pranchas de mapas.

